\documentclass[10pt]{article}
\usepackage[T1]{fontenc}
\usepackage[utf8]{inputenc}
\usepackage[french]{babel}
\usepackage{amsmath,amssymb,mathrsfs,mathpazo}
\usepackage{geometry}
\geometry{a4paper,margin=2cm}
\usepackage{graphicx,array,tabularx,tabularray,colortbl,multirow,multicol,booktabs,enumitem,pifont,ulem,fancybox,diagbox,cancel}
\usepackage{pstricks,pstricks-add,pst-plot,pst-tree,pst-node,pst-text,pst-func,pst-math,pst-bspline,pst-3dplot,pst-eucl,pst-solides3d}
\let\pscalculate\undefined
\usepackage{pst-bezier}
\usepackage{tikz,pgfplots,tkz-tab}
\usetikzlibrary{arrows,arrows.meta,patterns,calc,positioning,trees,decorations.markings,shapes.geometric,shapes.misc,shapes.arrows,intersections,decorations.pathmorphing,decorations.pathreplacing,backgrounds,shadows,babel}
\pgfplotsset{compat=1.18}
\usepgfplotslibrary{fillbetween}
\usepackage[np]{numprint}
\usepackage{listings,makecell,pas-tableur,twoopt}
\let\barre\undefined
\usepackage[pstricks]{tablvar}
\definecolor{darkgreen}{RGB}{0,100,0}
\providecommand{\red}{\color{red}}
\providecommand{\blue}{\color{blue}}
\providecommand{\green}{\color{green}}
\providecommand{\black}{\color{black}}
\providecommand{\white}{\color{white}}
\providecommand{\gray}{\color{gray}}
\providecommand{\cyan}{\color{cyan}}
\providecommand{\magenta}{\color{magenta}}
\providecommand{\yellow}{\color{yellow}}
\providecommand{\violet}{\color{violet}}
\providecommand{\orange}{\color{orange}}
\providecommand{\euro}{\texteuro}
\providecommand{\EUR}{\texteuro}
\providecommand{\ieme}{\textsuperscript{e}}
\providecommand{\degres}{\ensuremath{^\circ}}
\providecommand{\degree}{\ensuremath{^\circ}}
\providecommand{\vv}[1]{\overrightarrow{#1}}
\providecommand{\Coord}[1]{#1}
\providecommand{\decosix}{}
\providecommand{\starredbullet}{\ensuremath{\bullet}}
\pagestyle{empty}
\setlength{\parindent}{0pt}
\usepackage{adjustbox}\DeclareUnicodeCharacter{2236}{:}\DeclareUnicodeCharacter{2190}{\ensuremath{\leftarrow}}
\providecommand{\R}{\mathbb{R}}
\providecommand{\N}{\mathbb{N}}
\providecommand{\Z}{\mathbb{Z}}
\providecommand{\Q}{\mathbb{Q}}
\providecommand{\D}{\mathbb{D}}
\providecommand{\C}{\mathbb{C}}
\providecommand{\e}{{\mathrm{e}}}
\providecommand{\vect}[1]{\overrightarrow{#1}}
\providecommand{\vectt}[1]{\overrightarrow{\mathrm{#1}}}
\providecommand{\vvt}[1]{\overrightarrow{\mathrm{#1}}}
\providecommand{\np}[2][]{{#2~#1}}
\providecommand{\vv}[1]{{\overrightarrow{#1}}}
\renewcommand{\d}{{\mathrm{d}}}
\providecommand{\degres}{{^\circ}}
\providecommand{\degre}{{^\circ}}
\providecommand{\blue}{\color{blue}}
\providecommand{\red}{\color{red}}
\providecommand{\enskip}{\quad}
\providecommand{\llbracket}{\lbrack\!\lbrack}
\providecommand{\rrbracket}{\rbrack\!\rbrack}
\providecommand{\euro}{{\text{€}}}
\providecommand{\no}{{\text{n°}}}
\providecommand{\interval}[2]{{[#1;#2]}}
\providecommand{\intervFO}[2]{{[#1;#2[}}
\providecommand{\dotfill}{\ldots}

\providecommand{\ds}{\displaystyle}
\providecommand{\textding}[1]{#1}
\providecommand{\Oij}{\(\mathrm{O};\vec{\imath},\vec{\jmath}\)}
\providecommand{\Oijk}{\(\mathrm{O};\vec{\imath},\vec{\jmath},\vec{k}\)}
\providecommand{\Ouv}{\(\mathrm{O};\vec{u},\vec{v}\)}
\providecommand{\intervFF}[2]{\ensuremath{[#1;#2]}}
\renewcommand{\euro}{\ifmmode\text{\texteuro}\else\texteuro\fi}
\DeclareUnicodeCharacter{00B0}{\ifmmode^{\circ}\else\textdegree\fi}
\DeclareUnicodeCharacter{20AC}{\ifmmode\text{\texteuro}\else\texteuro\fi}
% Énoncés : Éducation nationale. Transcription et corrigés : APMEP, Denis Vergès et contributeurs. Découpage et classement : CoopMaths, Sébastien Lozano, Éric Elter et contributeurs. Adaptation de la présentation : Maths Explorer.
% CC BY-SA : https://coopmaths.fr/www/about/
% https://coopmaths.fr/alea/static/bac/2023/tex/bac_2023_03_sujet2_etrangers_groupe2_2_cor.tex
\begin{document}
\noindent\textit{APMEP / CoopMaths. CC BY-SA. Présentation adaptée par Maths Explorer.}\par\medskip
\medskip

\textbf{Partie A}

\medskip


\[f(x) = x - \ln (1 + x).\]

\begin{enumerate}
\item 
$\ln u$ est défini si $u > 0$, soit $1 + x > 0 \iff x > - 1$.

La fonction $f$ est définie sur l'intervalle $]-1~;~+ \infty[$.
\item 


En posant $u(x) = 1 + x$ et donc $u'(x) = 1$, on a alors  si $f'(x) = 1 - \frac{u'(x)}{u(x)} = 1 - \dfrac{1}{1 + x} = \dfrac{1 + x - 1}{1 + x} = \dfrac{x}{1 + x}$ sur $]-1~;~+ \infty[$.
\item 
	\begin{enumerate}
		\item 
On sait que $1 + x > 0$, donc le signe de la dérivée est celui de $x$ :
		
$\bullet~~f'(x) = 0 \iff x = 0$ ;

$\bullet~~f'(x) > 0 \iff x > 0$  : la fonction est croissante sur $[0~;~+ \infty[$ ;

$\bullet~~f'(x) < 0 \iff x < 0$ ; la fonction est décroissante sur $]-1~;~0[$.
		\item 
La fonction $f$ est décroissante sur $]-1~;~0[$ puis croissante sur $[0~;~+ \infty[$ : elle a donc un minimum en $x = 0, \: f(0 ) = 0 - 0 = 0$. 

Son minimum étant nul, on a donc sur $]-1~;~+ \infty[, f(x) \geqslant 0$.
	\end{enumerate}
\item 
	\begin{enumerate}
		\item 
On sait que quel que soit le réel $x$, $x = \ln \text{e}^x$. On peut donc écrire :


$f(x) = \ln \text{e}^x - \ln (1 + x) = \ln \left(\dfrac{\text{e}^x}{1 + x}\right)$ par propriété de la fonction logarithme népérien :

 $\ln a - \ln b = \ln \dfrac ab$


		\item 
En posant $u = 1 + x \iff x = u - 1$, on a $f(x) = \ln \left(\dfrac{\text{e}^x}{1 + x}\right) = \ln \left(\dfrac{\text{e}^{u - 1}}{u}\right) = \ln \left(\dfrac{\text{e}^{u}\times \text{e}^{-1}}{u}\right) = \ln \left(\dfrac{1}{\text{e}}\dfrac{\text{e}^{u}}{u}\right)$.
		
Or $x \to + \infty \implies x + 1 = u \to + \infty$ et on sait que $\displaystyle\lim_{u \to + \infty} \dfrac{\text{e}^u}{u} = + \infty$, donc $\displaystyle\lim_{u \to + \infty} \left(\dfrac{1}{\text{e}}\dfrac{\text{e}^{u}}{u}\right) = + \infty$ et enfin $\displaystyle\lim_{u \to + \infty} f(u) = + \infty$ ou $\displaystyle\lim_{x \to + \infty} f(x) = + \infty$.
	\end{enumerate}
\end{enumerate}

\medskip

\textbf{Partie B}

\medskip

On considère la suite $\left(u_n\right)$ définie par $u_0 = 10$ et, pour tout entier naturel $n$,

\[u_{n+1} = u_n - \ln \left(1 + u_n\right).\]



\medskip

\begin{enumerate}
\item Avec $n = 0$, la relation de récurrence donne :
$u_1 = u_0 - \ln (1 + u_0) = 10 - \ln (1 + 10) = 10 - \ln 11 \approx 7,602$.


\item 

\emph{Initialisation } : on a $u_0 = 10 \geqslant 0$ : l'inégalité est vraie au rang 0.

\emph{Hérédité } : Soit $n \in \N$ tel que $u_n \geqslant 0$ : d'après la question \textbf{3. a.} la fonction est croissante sur $[0~;~+ \infty[$, et $u_{n+1} = f\left(u_n\right) \geqslant 0$, donc $0 \leqslant u_{n+1}$ : l'inégalité est vraie au rang $n + 1$.

\emph{Conclusion } : l'inégalité est vraie au rang 0 et si elle est vraie au rang $n$, elle l'est aussi au rang $n + 1$ : d'après le principe de récurrence :

pour tout entier naturel $n$, on a $u_n \geqslant 0$.
\item 
On a $u_{n+1} = u_n - \ln \left(1 + u_n\right) \iff u_{n+1} - u_n = - \ln \left(1 + u_n\right)$.

Or on vient de démontrer que pour tout $n \in \N$, \: $u_n \geqslant 0 \implies 1 + u_n \geqslant 1$ puis par croissance de la fonction logarithme népérien $\ln (1 + u_n) \geqslant \ln 1$, soit  $\ln (1 + u_n) \geqslant 0 \iff - \ln (1 + u_n) \leqslant 0$.

Donc $u_{n+1} - u_n \leqslant 0$ montre que la suite $\left(u_n\right)$ est décroissante.
\item 

La suite $\left(u_n\right)$ étant décroissante et minorée par $0$ converge vers une limite $\ell$ telle que $\ell \geqslant 0$
\item 
On a $u_{n+1} = f\left(u_n\right)$. On a vu que cette fonction est dérivable donc continue sur $]- 1~;~+ \infty[$ ; de plus la suite $\left(u_n\right)$ converge vers $\ell$, donc à la limite :

$\ell = \ell - \ln \left(1 + \ell\right) \iff \ln \left(1 + \ell\right) = 0$ ou $\ln \left(1 + \ell\right) = \ln 1 \iff 1 + \ell = 1 \iff \ell = 0$.

Conclusion : $\displaystyle\lim_{n \to + \infty} u_n = 0$.
\end{enumerate}

\medskip

\medskip

\end{document}
